\subsection{Validation post-deploiement et resultats}
\label{sec:validation-postdeploy}

\subsubsection{La phase post-déploiement et l'amélioration continue}
\label{subsec:post-deployment-continuous-improvement}

Traditionnellement, la sécurité dans le cycle de développement logiciel (SSDLC) se focalise principalement sur la phase pré-déploiement (SAST, SCA, scans de secrets, builds sécurisés). Cependant, considérer que le passage en production marque la fin des vérifications de sécurité constitue une faille méthodologique majeure. Une fois l'application déployée au sein d'un cluster Kubernetes, de nouveaux vecteurs d'attaques dynamiques, des dérives de configuration (\textit{configuration drift}) ou des vulnérabilités de type \textit{zero-day} peuvent survenir à tout moment.

\begin{quote}
  \textit{« Dans SecureRAG Hub, le déploiement Kubernetes n'est pas la fin de la chaîne DevSecOps. Il constitue le point de départ d'une phase de validation post-déploiement visant à vérifier le bon fonctionnement applicatif, la conformité des images exécutées, le respect des politiques de sécurité Kubernetes, la disponibilité des preuves techniques et l'alimentation d'une boucle d'amélioration continue. »}
\end{quote}

Les composants clés qui structurent cette phase post-déploiement sont les suivants :

\begin{itemize}
  \item \textbf{La validation post-déploiement} : Elle agit comme une barrière d'assurance qualité et de sécurité en continu. Elle permet de s'assurer, de manière déterministe, que les descripteurs Kubernetes ont été appliqués avec succès, que les conteneurs sont en état de marche (\texttt{Running} et \texttt{Ready}) et qu'aucun service critique n'est indisponible en raison de dysfonctionnements de démarrage (comme les états \texttt{CrashLoopBackOff} ou \texttt{ImagePullBackOff}).
  \item \textbf{La preuve runtime (\textit{Runtime Image Proof})} : Dans une chaîne logistique logicielle sécurisée (\textit{Supply Chain Security}), il est crucial de garantir l'authenticité des charges de travail exécutées. La preuve runtime consiste à interroger directement l'API Kubernetes pour obtenir l'identifiant unique (SHA256 digest) de chaque conteneur en cours d'exécution. Cette valeur est ensuite confrontée aux digests officiels et signés lors de la phase de CI/CD. Cela élimine définitivement les attaques basées sur l'usurpation d'images ou les déploiements non autorisés d'images tierces utilisant des tags mutables comme \texttt{latest}.
  \item \textbf{La sécurité runtime (\textit{Kernel \& Calls System Audit})} : Les contrôles statiques au niveau de l'admission (comme Kyverno) ou les configurations de pods (\texttt{SecurityContext}) restreignent les privilèges de base mais ne bloquent pas toutes les tentatives d'intrusion actives. La sécurité runtime s'appuie sur des outils d'analyse système avancés comme \textbf{Falco}. En écoutant directement les appels système au niveau du noyau Linux (via eBPF), le système est capable de détecter instantanément et d'alerter lors de comportements suspects, tels que le lancement d'un terminal interactif dans un conteneur web ou une tentative d'écriture dans un répertoire système en lecture seule.
  \item \textbf{L'observabilité} : L'observabilité moderne ne sert pas uniquement à surveiller les pannes matérielles ou les temps de réponse. En DevSecOps, elle corrèle les données de performance applicative (SLO de disponibilité et de latence) avec les métriques de sécurité (violations de politiques Kyverno, alertes Falco, logs système Loki). Cette centralisation offre aux équipes de détection une visibilité à 360 degrés sur la santé globale et le niveau d'exposition du cluster.
  \item \textbf{Le support pack (\textit{Evidence Collection})} : Dans le cadre d'audits de conformité de sécurité (comme ISO 27001 ou SOC2) ou d'une soutenance de fin d'études, il est indispensable de pouvoir présenter des preuves irréfutables de la conformité du système. Le support pack regroupe de manière automatique, à chaque déploiement, l'ensemble des rapports générés (rapports de scans Semgrep/Trivy, manifestes appliqués, digests réels vérifiés par Cosign, logs de validation) sous une archive compressée et signée.
  \item \textbf{La boucle de feedback fermée (\textit{Feedback Loop})} : Le DevSecOps s'inscrit dans un modèle d'amélioration continue. Tout dysfonctionnement détecté au runtime (que ce soit une faille applicative détectée par smoke test ou une violation de règle système Falco) doit obligatoirement remonter sous forme d'alerte, être analysé par l'équipe d'ingénierie, et donner lieu à un correctif (correction de code Laravel, durcissement de configuration Kustomize, ou ajustement de politique Kyverno). Cette boucle fermée de feedback boucle l'ensemble du cycle de vie du logiciel et renforce de manière continue la résilience du système.
\end{itemize}

\subsubsection{Objectifs et méthodologie de validation}
\label{subsubsec:objectifs-methodologie-validation}

Cette section presente la validation post-deploiement de la plateforme \textit{SecureRAG Hub} sur un cluster Kubernetes local fonde sur \texttt{kind}. L'objectif n'est pas uniquement de verifier que les conteneurs demarrent, mais aussi de demontrer qu'une chaine DevSecOps reproductible peut etre executee localement, que les controles de securite essentiels sont actifs, et qu'un jeu de preuves reutilisable est produit automatiquement apres deploiement.

\subsubsection{Objectifs de validation}
\label{subsubsec:objectifs-validation-detail}

La validation post-deploiement a ete structuree autour de quatre objectifs complementaires :

\begin{itemize}
  \item verifier le demarrage du cluster local, du registre local et des workloads Kubernetes ;
  \item confirmer l'accessibilite des workloads Laravel officiels ;
  \item exercer un premier niveau de validation securite sur les \textit{security contexts}, les \textit{NetworkPolicy} et la disponibilite du composant \texttt{audit-security-service} ;
  \item produire un rapport synthetique exploitable dans le memoire et en soutenance.
\end{itemize}

Mise a jour du perimetre : la validation officielle actuelle porte sur \texttt{portal-web}, \texttt{auth-users}, \texttt{chatbot-manager}, \texttt{conversation-service} et \texttt{audit-security-service}. Les anciens composants Python/RAG et Ollama sont conserves comme historique technique, mais exclus de la build/deploy officielle tant que leurs sources applicatives ne sont pas restaurees.

\subsubsection{Protocole experimental}

Le premier essai reel a ete execute le 7 avril 2026 sur une machine locale \texttt{arm64}. Le protocole retenu suit une sequence simple et reproductible :

\begin{enumerate}
  \item creation du cluster \texttt{kind} et du registre local ;
  \item construction locale des images OCI des cinq workloads officiels ;
  \item deploiement des manifests Kubernetes via Kustomize sur l'overlay \texttt{dev} ;
  \item execution des scripts de validation post-deploiement ;
  \item generation d'un rapport global de synthese.
\end{enumerate}

Les commandes principales utilisees sont les suivantes :

\begin{verbatim}
bash scripts/deploy/create-kind.sh
bash scripts/deploy/build-local-images.sh
bash scripts/deploy/deploy-kind.sh

bash scripts/validate/smoke-tests.sh
bash scripts/validate/security-smoke.sh
bash scripts/validate/e2e-functional-flow.sh
bash scripts/validate/security-adversarial-advanced.sh
bash scripts/validate/generate-validation-report.sh
\end{verbatim}

Le deploiement est realise dans l'espace de noms \texttt{securerag-hub}. Les images applicatives sont construites localement puis poussees dans le registre \texttt{localhost:5001}. L'acces externe de demonstration est expose via le \texttt{NodePort} du portail sur \texttt{http://localhost:8081/health}.

\subsubsection{Jeu de validations mis en oeuvre}

Le dispositif de validation repose sur plusieurs scripts complementaires.

\paragraph{Validation de base et controles Kubernetes.}
Le script \texttt{security-smoke.sh} controle la presence des \textit{ServiceAccount} dedies, l'activation des \textit{NetworkPolicy} et les principaux parametres de durcissement des conteneurs, en particulier \texttt{runAsNonRoot}, \texttt{readOnlyRootFilesystem} et \texttt{allowPrivilegeEscalation=false} lorsque ces contraintes sont applicables.

\paragraph{Validation fonctionnelle de bout en bout.}
Le script \texttt{e2e-functional-flow.sh} verifie l'existence des services Kubernetes officiels, l'accessibilite des endpoints internes \texttt{/health} depuis un pod de validation, puis l'exposition externe du portail via \texttt{http://localhost:8081/health}.

\paragraph{Validation RAG legacy.}
Le script \texttt{rag-smoke.sh} ne fait plus partie de la validation officielle par defaut. Il retourne un statut \texttt{PRET\_NON\_EXECUTE} sauf activation explicite du perimetre legacy, car les sources Python correspondantes ne sont plus presentes.

\paragraph{Validation adversariale.}
Le script \texttt{security-adversarial-advanced.sh} verifie la disponibilite du composant \texttt{audit-security-service} et controle que l'endpoint d'audit refuse les appels non autorises.

\subsubsection{Mode d'interpretation des resultats}

Les scripts de validation utilisent trois statuts :

\begin{itemize}
  \item \texttt{PASS} : le comportement attendu est observe ;
  \item \texttt{FAIL} : une anomalie bloquante est detectee ;
  \item \texttt{SKIP} : le test n'est pas applicable ou l'endpoint metier n'est pas encore implemente.
\end{itemize}

Ce choix est adapte au contexte d'une preuve de concept. Il permet de demontrer un socle DevSecOps complet sans supposer que tous les endpoints metier definitifs sont deja finalises.

\subsubsection{Resultats obtenus lors du premier run reel}

Le premier run reel historique a permis d'obtenir les observations suivantes. Ces observations ne representent plus le perimetre officiel actuel, mais expliquent pourquoi la trajectoire a ete recentree sur le runtime Laravel-first :

\begin{itemize}
  \item le cluster \texttt{securerag-dev} a ete cree avec succes ;
  \item le registre local \texttt{localhost:5001} a ete initialise correctement ;
  \item les anciennes images applicatives ont ete construites et poussees localement ;
  \item les microservices \texttt{api-gateway}, \texttt{auth-users}, \texttt{chatbot-manager}, \texttt{llm-orchestrator}, \texttt{security-auditor} et \texttt{knowledge-hub} ont atteint l'etat \texttt{Running} ;
  \item le composant \texttt{qdrant} a ete deploye avec succes sous forme de \texttt{StatefulSet} ;
  \item le composant \texttt{ollama} est reste bloque a l'etat \texttt{ContainerCreating} lors du premier essai ;
  \item l'endpoint externe \texttt{http://localhost:8080/healthz} a repondu correctement.
\end{itemize}

Le rapport synthetique genere automatiquement dans \texttt{artifacts/validation/validation-summary.md} indique le bilan suivant :

\begin{center}
\begin{tabular}{|l|c|}
\hline
Statut & Nombre \\
\hline
\texttt{PASS} & 15 \\
\texttt{FAIL} & 1 \\
\texttt{SKIP} & 2 \\
\hline
\end{tabular}
\end{center}

Le detail des principaux tests est resume ci-dessous :

\begin{itemize}
  \item \textbf{Validation fonctionnelle E2E :} succes. Les services Kubernetes existent, les endpoints internes \texttt{/healthz} sont joignables depuis un pod de validation, l'acces externe via \texttt{NodePort} fonctionne, et la racine de \texttt{api-gateway} repond ;
  \item \textbf{Validation securite :} succes partiel. Les \textit{ServiceAccount} dedies et les \textit{NetworkPolicy} sont bien presents. Le composant \texttt{security-auditor} est joignable. L'endpoint \texttt{/analyze} n'etant pas encore implemente, le test correspondant est en \texttt{SKIP} ;
  \item \textbf{Validation RAG :} echec. Le test \texttt{rag-smoke.sh} echoue car \texttt{ollama} n'est pas encore joignable, ce qui empeche de considerer la chaine RAG comme entierement operationnelle.
\end{itemize}

\subsubsection{Analyse technique des resultats}

Le socle DevSecOps de la plateforme est donc valide historiquement sur les points suivants, et doit etre revalide sur le perimetre Laravel-first avec les commandes courantes :

\begin{itemize}
  \item la creation du cluster local et du registre local est reproductible ;
  \item la chaine de build locale des microservices applicatifs est fonctionnelle ;
  \item l'overlay Kustomize \texttt{dev} deploie correctement les ressources Kubernetes ;
  \item l'exposition externe via \texttt{NodePort} est operationnelle dans le scenario teste ;
  \item les politiques reseau et les controles de securite minimum sont bien appliques ;
  \item les scripts de validation produisent des artefacts exploitables automatiquement.
\end{itemize}

L'anomalie principale identifiee historiquement concernait \texttt{ollama}. Cette anomalie n'est plus un bloqueur du scenario officiel, car le perimetre legacy RAG/Ollama est exclu de la build/deploy officielle tant que ses sources et preuves ne sont pas restaurees.

Il convient egalement de noter un point de vigilance securite : le namespace applique les labels \texttt{pod-security} avec \texttt{enforce=baseline} et \texttt{audit/warn=restricted}. Dans ce cadre, \texttt{ollama} remonte un avertissement sur \texttt{runAsNonRoot}, ce qui est acceptable pour le bootstrap local, mais devra etre traite dans une iteration ulterieure si l'objectif est d'approcher plus strictement le profil \texttt{restricted}.

\subsubsection{Artefacts de preuve produits}

Le protocole a genere automatiquement plusieurs fichiers dans le repertoire \texttt{artifacts/validation/}, notamment :

\begin{itemize}
  \item \texttt{validation-summary.md} ;
  \item \texttt{e2e-functional-flow.txt} ;
  \item \texttt{rag-smoke.txt} ;
  \item \texttt{security-adversarial.txt} ;
  \item \texttt{k8s-get-all.txt} ;
  \item \texttt{k8s-pods.txt} ;
  \item \texttt{k8s-pvc.txt} ;
  \item \texttt{k8s-networkpolicy.txt} ;
  \item \texttt{portal-web-describe.txt} ;
  \item \texttt{auth-users-describe.txt} ;
  \item \texttt{chatbot-manager-describe.txt} ;
  \item \texttt{conversation-service-describe.txt} ;
  \item \texttt{audit-security-service-describe.txt}.
\end{itemize}

Ces fichiers constituent un jeu de preuves directement reutilisable en soutenance : etat du cluster, ressources deployees, stockage persistant, politiques reseau, validation fonctionnelle, verification securite et synthese finale.

\subsubsection{Corrections et stabilisation apres le premier run}

Le premier run reel a joue un role utile de diagnostic. Il a mis en evidence plusieurs points qui ont ensuite ete corriges pour fiabiliser l'exploitation locale.

\paragraph{Rationalisation du runtime officiel.}
La correction principale a consiste a retirer les composants Python/RAG non exploitables de la build/deploy officielle et a integrer les services Laravel réellement presents : \texttt{conversation-service} et \texttt{audit-security-service}.

\paragraph{Activation des addons de cluster.}
Dans un second temps, \texttt{metrics-server} et \texttt{Kyverno} sont prepares par manifests et runbooks, mais leur validation reste dependante d'un cluster actif. Une preuve ne doit etre declaree \texttt{TERMINEE} qu'apres execution de :

\begin{verbatim}
kubectl top pods -n securerag-hub
kubectl get hpa -n securerag-hub
kubectl get clusterpolicy
kubectl get policyreport -A
\end{verbatim}

Sans ces commandes et leurs sorties archivees, le statut correct reste \texttt{DEPENDANT\_DE\_L\_ENVIRONNEMENT}.

\paragraph{Correction des pods de validation sous policy Cosign.}
Une fois Kyverno actif, un point de vigilance supplementaire est apparu : les pods ephemeres de validation bases sur une image \texttt{securerag-hub-*} du registre local pouvaient etre controles par la policy de verification Cosign. Lorsque le registre local n'etait pas resolu de la meme maniere depuis le cluster, l'admission pouvait etre refusee.

La correction retenue a ete d'utiliser une image publique de validation pour les pods temporaires de test, afin de separer clairement :

\begin{itemize}
  \item la verification de la sante fonctionnelle du cluster ;
  \item la verification de la chaine de confiance des images applicatives.
\end{itemize}

Cette distinction renforce la lisibilite de la validation post-deploiement et evite de confondre un probleme de registre local avec un probleme d'accessibilite des services.

\subsubsection{Interpretation pour la soutenance}

Les resultats obtenus permettent de soutenir les affirmations suivantes :

\begin{itemize}
  \item la plateforme n'est pas seulement decrite ou conteneurisee ; elle est effectivement deployable localement ;
  \item la chaine DevSecOps est operationnelle de la construction des images jusqu'a la validation post-deploiement ;
  \item l'architecture microservices Laravel est correctement exposee et testable ;
  \item les controles de securite Kubernetes et les tests adversariaux sont integres des le socle technique ;
  \item les points restant a prouver concernent l'environnement runtime : cluster actif, metrics-server, Kyverno, registry, Syft et Cosign.
\end{itemize}

\subsubsection{Conclusion}

La validation post-deploiement de \textit{SecureRAG Hub} confirme la solidite du socle DevSecOps du projet, a condition de distinguer strictement les preuves deja executees des preuves dependantes de l'environnement. Le runtime officiel est desormais Laravel-first ; la chaine legacy RAG/Ollama reste hors perimetre officiel tant qu'elle n'est pas restauree et revalidee.
